\chapter{Downstream Tasks}\label{ch:downstream}
Four tasks with objective metrics and matched baselines test whether the
per-note posterior is useful beyond imputation, and all outcomes are
reported. Table~\ref{tab:downstream} summarises them, and the results of the
two-stage pipeline during development stand beside them in the development
record.\footnote{\texttt{docs/downstream\_tasks\_results.md} in the
repository.} Every task was run under the two-stage pipeline during
development and evaluated again under the proposed model with identical
corruptions and masks wherever the task permits pairing. Two
whole-performance summary tasks run during development are recorded there
and not reported here: musical-era classification, a negative result under
both models, and performer identification on the Vienna 4x22 corpus, mixed
under both, because smoothing removes the idiosyncrasy that identifies a
pianist. They are the evidence for the second half of the structure boundary
of \S\ref{sec:boundary}.

\begin{table}[h]\centering
\caption{Downstream tasks under the proposed model. The results of the
two-stage pipeline during development are recorded in the development
record.}
\label{tab:downstream}
\footnotesize
\begin{tabular}{@{}p{4.6cm}p{8.2cm}@{}}
\toprule
Task & Outcome under the proposed model \\
\midrule
Error spotting (anomaly) & best on all channels ($v$ 0.995) \\
Denoising, known noise & better (0.182/0.276, cov 0.90) \\
Denoising, blind noise & no improvement in either parameterization \\
Completion from excerpt & partial, with an extrapolation boundary (\S\ref{sec:boundary}) \\
Selective prediction & better on the confident half, transfers (confident-half RMSE 0.076) \\
\bottomrule
\end{tabular}
\end{table}

\section{Error spotting}
Error spotting corrupts a small fraction of notes and ranks every note by
its suspiciousness, so the measure is the quality of the ranking. The
two-stage graph model ranks better than every structure-free baseline on
every channel, with AUROC rising from $0.94 \to 0.98$ at $3\sigma$ and from
$0.85 \to 0.94$ at $2\sigma$. The proposed model, which scores each note by
its joint leave-one-out predictive surprise, improves further on all three
channels, reaching a velocity AUROC of $0.995$, an articulation AUROC of
$0.986$ and a timing AUROC of $0.978$. Error spotting is therefore the task
on which structure helps most clearly, and the advantage is larger under the
proposed model than under the two-stage pipeline.

\section{Denoising}
Denoising recovers the true values of the per-note variables from added
noise. With the noise level known, structure is better than per-note
shrinkage on both error and calibration, with all 8 paired differences having
intervals that exclude zero and coverage near $0.90$. The proposed model's
oracle-noise form is slightly better still, with RMSE $0.182$ and $0.276$ at
the two noise levels against $0.186$ and $0.282$, and coverage $0.90$. A
negative result holds unchanged under both models: estimating the noise level
blind from a single piece is unreliable in \emph{both} parameterizations,
because the intervals collapse, so transcription noise must be treated as
known for each system.

\section{Completion from an excerpt}
Completion predicts the expression of the rest of a piece from its opening.
Interpolation, with scattered hidden notes, is where the model is clearly
better, with RMSE falling from $0.450 \to 0.393$ and NLL from $0.21 \to 0.10$
at half the notes observed. When the prediction extrapolates far past the
excerpt, the music-model prediction carries it and the graph is neutral or
slightly harmful. The proposed model makes this boundary sharp. Its per-piece
feature weights, fit on the opening excerpt alone, can extrapolate wildly,
and one evaluation with a prefix of 25\% reached an RMSE of
$1.6\times 10^4$, so the cross-piece read-out head, which never adapts to the
piece, is the appropriate tool for excerpt extrapolation. Per-piece Bayesian
adaptation needs observed coverage of the feature space, and rendering from a
short excerpt does not provide it.

\section{Selective prediction}
Selective prediction discards the least-confident predictions and measures
the error on those that remain. Discarding lowers error most clearly for
articulation, where keeping the confident half cuts its RMSE from $0.83 \to
0.54$, and for catching the model's own misfires, where a timing blow-up at
$0.61$ overall falls to $0.07$ on the confident half. Only the graph provides
the per-note confidence that this triage needs, because the no-graph
baselines rank at piece level only. Under the proposed model the triage is
slightly better still, with a confident-half RMSE of $0.076$ against $0.080$
at matched abstention, so the property transfers from the two-stage pipeline
to the proposed model.

\section{Two boundaries}\label{sec:boundary}
Across the tasks, two boundaries emerge. Boundary (1), the structure
boundary, holds under both models: the graph prior helps whenever notes are
judged one by one for accuracy or confidence, and does not help when a
performance is collapsed into a summary. Boundary (2), the adaptation
boundary, belongs to the proposed model: per-piece Bayesian feature weighting
is better on every interpolation-style task and breaks down at excerpt
extrapolation, where the cross-piece head remains the appropriate tool. Both
boundaries are statements this thesis makes, and each is bounded by evidence
on both sides.

% =====================================================================
